\chapter{Settings}
\label{ch:settings}

Settings are kept on this computer and apply at once. The list on the left
groups them as on the Tools page: \emph{Modbench} (About, Updates),
\emph{Display} (Appearance, Values) and \emph{Devices} (Connection, Device
library).

\section{About}

The version of Modbench and whether it is up to date, with \ui{Check for
updates}. Below it are links to this manual, the release notes, a form to
report a problem, the source code and the project's Ko-fi page.

\screenshot[8cm]{about}{The About section.}
  {Settings on About: the app, its version and the links to help and to the
   project.}

\screenshot[8cm]{settings}{The Settings page, on Appearance.}
  {The Settings page: the list of sections on the left, and Appearance open
   with the three theme previews and the text size.}

\section{Appearance}

\begin{description}[style=nextline, leftmargin=0pt, font=\intermedium]
  \item[Theme] Light, dark, or as the system.
  \item[Text size] Smaller to see more at once, larger to read from further
    away. \ui{Reset} goes back to 100\,\%.
\end{description}

\section{Values}

\begin{description}[style=nextline, leftmargin=0pt, font=\intermedium]
  \item[Highlight changed values] How long a value that changed stays
    highlighted: off, or 1 to 30 seconds.
\end{description}

\section{Device library}

\begin{description}[style=nextline, leftmargin=0pt, font=\intermedium]
  \item[Library folder] Where your maps are kept. \ui{Open folder} shows it.
  \item[Shared maps] Links to the map format and to the repository of maps
    shared by other users.
\end{description}

\section{Updates}

\begin{description}[style=nextline, leftmargin=0pt, font=\intermedium]
  \item[Check for updates when Modbench starts] On by default. Installs from
    the APT repository update with \val{sudo apt upgrade} instead.
\end{description}

\section{Connection}

\begin{description}[style=nextline, leftmargin=0pt, font=\intermedium]
  \item[Remember the last connection] Port, baud rate, address, device and
    timing are filled in again when Modbench starts.
  \item[Remember register choices for each device] Which registers are read,
    the byte order and Loop come back for each device (map, connection and
    unit ID). Restoring them never writes or starts a test.
  \item[Forget saved device choices] Clears what was remembered per device.
  \item[Confirm before writing to a device] Asks before every write.
  \item[Reconnect when the connection is lost] While a test runs, a connection
    that fails is opened again, and the test goes on.
  \item[Read-only mode] Blocks every write, also on a connection that is
    already open.
  \item[Write single registers with FC16] One holding register is written
    with FC06 (Write Single Register) unless this is on; then FC16 (Write
    Multiple Registers) is used, for devices that only accept FC16.
  \item[CSV separator] For saved logs, exported values and recordings. Use a
    semicolon for spreadsheets set to a language with a decimal comma.
  \item[Requests kept in the log] Older requests are dropped first.
\end{description}

\section{What Modbench keeps}

Everything stays on this computer: settings, the last connection, choices per
device, alarms, templates and the library folder. The only network requests
besides your devices are downloads: the update check, the online library, and
the vendors' icons from their websites.
